%! The Quadratic Formula & the Discriminant — Unit 2, Lesson 2.5 — Homework (Answer Key)
\documentclass[10pt]{article}
\usepackage{algebra2-article}
\usepackage{algebra2-key}   % pulls in -boxes; do NOT also load -boxes
\newcommand{\TallMath}[1]{$\displaystyle #1\rule[-1.4em]{0pt}{3.2em}$}

\begin{document}

\pageheader{Unit 2, Lesson 2.5}{Homework --- Answer Key}

\vspace{0.10in}

\begin{notesbox}{Practice}
Show your work. Write in standard form first, read $a,b,c$ with signs, and simplify every answer (use
$a+bi$ when the roots are complex).

\begin{enumerate}[leftmargin=1.9em, itemsep=9pt]
  \item \textbf{Standard form, then $a,b,c$.}
    \begin{enumerate}[label=(\alph*), leftmargin=1.9em, itemsep=4pt]
      \item $x^2=6x-2 \Rightarrow$ \ans{$x^2-6x+2=0$}\quad $a=\ans{$1$}\ b=\ans{$-6$}\ c=\ans{$2$}$
      \item $5x-3=2x^2 \Rightarrow$ \ans{$2x^2-5x+3=0$}\quad $a=\ans{$2$}\ b=\ans{$-5$}\ c=\ans{$3$}$
    \end{enumerate}

  \item \textbf{Discriminant \& classify} (compute $b^2-4ac$; state two real / one repeated / two complex).
    \[ x^2-2x-8=0:\ \ans{$36$; two real} \qquad x^2+x+1=0:\ \ans{$-3$; two complex} \]
    \[ 9x^2-12x+4=0:\ \ans{$0$; one repeated} \qquad x^2+5x+2=0:\ \ans{$17$; two real} \]

  \item \textbf{Solve} (real, rational).
    \begin{enumerate}[label=(\alph*), leftmargin=1.9em, itemsep=4pt]
      \item $x^2-x-6=0$
        \begin{work}
          x &= \frac{1\pm\sqrt{(-1)^2-4(1)(-6)}}{2(1)} \\
            &= \frac{1\pm\sqrt{25}}{2} \\
            &= \frac{1\pm5}{2} \\
            &= 3,\ -2
        \end{work}
      \item $2x^2-7x+3=0$
        \begin{work}
          x &= \frac{7\pm\sqrt{(-7)^2-4(2)(3)}}{2(2)} \\
            &= \frac{7\pm\sqrt{25}}{4} \\
            &= \frac{7\pm5}{4} \\
            &= 3,\ \frac{1}{2}
        \end{work}
    \end{enumerate}

  \item \textbf{Solve} (real, irrational --- simplest radical form).
    \begin{enumerate}[label=(\alph*), leftmargin=1.9em, itemsep=4pt]
      \item $x^2-6x+7=0$
        \begin{work}
          x &= \frac{6\pm\sqrt{(-6)^2-4(1)(7)}}{2(1)} \\
            &= \frac{6\pm\sqrt{8}}{2} \\
            &= \frac{6\pm2\sqrt{2}}{2} \\
            &= 3\pm\sqrt{2}
        \end{work}
      \item $x^2+4x+1=0$
        \begin{work}
          x &= \frac{-4\pm\sqrt{4^2-4(1)(1)}}{2(1)} \\
            &= \frac{-4\pm\sqrt{12}}{2} \\
            &= \frac{-4\pm2\sqrt{3}}{2} \\
            &= -2\pm\sqrt{3}
        \end{work}
    \end{enumerate}

  \item \textbf{Solve} (complex --- $a+bi$ form).
    \begin{enumerate}[label=(\alph*), leftmargin=1.9em, itemsep=4pt]
      \item $x^2-2x+10=0$
        \begin{work}
          x &= \frac{2\pm\sqrt{(-2)^2-4(1)(10)}}{2(1)} \\
            &= \frac{2\pm\sqrt{-36}}{2} \\
            &= \frac{2\pm6i}{2} \\
            &= 1\pm3i
        \end{work}
      \item $x^2+16=0$
        \begin{work}
          x &= \frac{0\pm\sqrt{0^2-4(1)(16)}}{2(1)} \\
            &= \frac{\pm\sqrt{-64}}{2} \\
            &= \frac{\pm8i}{2} \\
            &= \pm4i
        \end{work}
      \item $2x^2+2x+1=0$
        \begin{work}
          x &= \frac{-2\pm\sqrt{2^2-4(2)(1)}}{2(2)} \\
            &= \frac{-2\pm\sqrt{-4}}{4} \\
            &= \frac{-2\pm2i}{4} \\
            &= -\frac{1}{2}\pm\frac{1}{2}i
        \end{work}
    \end{enumerate}

  \item \textbf{Discriminant $\to$ graph.} State how many times each parabola crosses the $x$-axis.
    \[ b^2-4ac=49:\ \ans{$2$} \qquad b^2-4ac=0:\ \ans{$1$} \qquad b^2-4ac=-4:\ \ans{$0$} \]

  \item \textbf{Choose the most efficient method} (square roots / factoring / quadratic formula --- you need not solve).
    \[ x^2-25=0:\ \ans{square roots} \quad x^2-7x+12=0:\ \ans{factoring} \quad x^2+3x+5=0:\ \ans{formula} \]

  \item \textbf{Explain.} A classmate says ``the discriminant is negative, so the equation has \emph{no}
        solutions.'' What is the correct statement?
        \ansline{A negative discriminant means no \emph{real} solutions --- but there are still \emph{two} solutions: a complex-conjugate pair $a\pm bi$. Every quadratic has solutions over the complex numbers.}
\end{enumerate}
\end{notesbox}

\begin{extensionbox}
\textbf{1. Model.} A firework is launched so its height is $h(t)=-16t^2+32t+48$ (feet, $t$ in seconds).
Solve $h(t)=0$ with the formula to find when it lands; reject the impossible root.
\begin{work}
  t &= \frac{-32\pm\sqrt{32^2-4(-16)(48)}}{2(-16)} \\
    &= \frac{-32\pm\sqrt{4096}}{-32} \\
    &= \frac{-32\pm64}{-32} \\
    &= -1,\ 3 \quad\text{reject } t=-1,\ \text{so it lands at } t=3\text{ s}
\end{work}

\textbf{2. Parameter.} For what value of $k$ does $x^2+6x+k=0$ have exactly one repeated real root?
\ans{$k=9$}\; For what values of $k$ does it have two complex roots? \ans{$k>9$}

\textbf{3. Verify a complex root.} Substitute $x=1+3i$ into $x^2-2x+10$ (from \#5a) and show it gives $0$.
\begin{work}
  (1+3i)^2-2(1+3i)+10 &= (1+6i+9i^2)-(2+6i)+10 \\
                      &= (1+6i-9)-(2+6i)+10 \\
                      &= (-8+6i)-(2+6i)+10 \\
                      &= 0
\end{work}
\end{extensionbox}

\begin{spiralbox}
\textbf{Coming next --- Lesson 2.6: Systems Involving Quadratics.} You can now solve \emph{any} single
quadratic. Next you'll solve \emph{systems} --- a line and a parabola, or two parabolas --- by substitution.
The catch: substitution collapses the system into one quadratic equation, which you'll finish with today's
formula, and the discriminant will tell you how many times the two graphs intersect.
\end{spiralbox}

\end{document}
